\chapter{Two canonical model systems}
\label{ch:model-systems}

\section*{Chapter overview}

This chapter solves two systems that, between them, model most of
the dynamics encountered in circuit QED: the quantum harmonic
oscillator and the two-level system (TLS). They are
\emph{idealisations}, not faithful descriptions of any physical
hardware --- but they are the building blocks out of which the more
realistic models in later chapters are assembled.

LC resonators, microwave cavities and transmission-line modes are,
in the lossless and linear limit, exact harmonic oscillators.
Superconducting \emph{qubits}, by contrast, are not strictly
two-level: they are weakly anharmonic oscillators with an infinite
ladder of energy levels, designed so that the
$\ket{0}\!\to\!\ket{1}$ transition frequency differs slightly from
$\ket{1}\!\to\!\ket{2}$. The TLS is the standard idealisation that
ignores all but the lowest two levels; it is accurate when drives
and couplings are slow and weak enough that population of $\ket{2}$
and beyond stays negligible. When such leakage matters --- for
instance during fast gates or strong drives --- the more honest
picture is the (anharmonic) oscillator, and one has to keep track
of the full ladder.

We will routinely move between the two. A transmon (\cref{ch:circuit-qed})
behaves as a TLS for slow control and as an anharmonic oscillator
when leakage is on the agenda. The toolbox developed in this
chapter --- ladder operators, the Pauli algebra, the Bloch sphere,
and Rabi oscillations --- is what we need in either limit. The
chapter ends by combining a TLS with a coherent drive, producing
the Rabi oscillations that underlie all single-qubit control.

% --------------------------------------------------------------
\section{The quantum harmonic oscillator}
\label{sec:qho}

\begin{phenomenon}
A particle in a quadratic potential, an LC circuit, and a single
mode of the electromagnetic field all share the same equation of
motion: simple harmonic motion at a single frequency $\omega$. Their
quantum versions therefore share a single Hilbert space structure
and a single energy spectrum.
\end{phenomenon}

\begin{statement}[Harmonic-oscillator spectrum]
\label{stmt:qho-spectrum}
The Hamiltonian
\begin{equation}
\op{H} \;=\; \frac{\op{p}^2}{2m} + \tfrac12 m\omega^2 \op{x}^2
\end{equation}
has eigenvalues $E_n=\hbar\omega(n+\tfrac12)$ for $n=0,1,2,\ldots$,
with non-degenerate eigenstates $\ket{n}$ that span the Hilbert
space.
\end{statement}

\paragraph{Derivation via ladder operators.}
Define the dimensionless ladder operators
\begin{equation}
\label{eq:ladder-def}
\hat{a}   = \sqrt{\frac{m\omega}{2\hbar}}\Bigl(\op{x}+\frac{i}{m\omega}\op{p}\Bigr),
\qquad
\hat{a}^\dagger = \sqrt{\frac{m\omega}{2\hbar}}\Bigl(\op{x}-\frac{i}{m\omega}\op{p}\Bigr).
\end{equation}
Using $[\op{x},\op{p}]=i\hbar\,\id$ one immediately obtains
\begin{equation}
[\hat{a},\hat{a}^\dagger] = \id,\qquad
\op{H} = \hbar\omega\bigl(\hat{a}^\dagger \hat{a} + \tfrac12\bigr).
\end{equation}
The number operator $\op{N}=\hat{a}^\dagger\hat{a}$ commutes with $\op{H}$, so
they share an eigenbasis. From $[\op{N},\hat{a}]=-\hat{a}$ and
$[\op{N},\hat{a}^\dagger]=+\hat{a}^\dagger$
one finds that $\hat{a}$ and $\hat{a}^\dagger$ shift the eigenvalue of
$\op{N}$ by $\mp 1$:
\begin{equation}
\label{eq:ladder-action}
\hat{a}\ket{n} = \sqrt{n}\,\ket{n-1},
\qquad
\hat{a}^\dagger\ket{n} = \sqrt{n+1}\,\ket{n+1}.
\end{equation}
Positivity, $\bra{\psi}\hat{a}^\dagger \hat{a}\ket{\psi}\geq0$, forces $n\geq0$;
a unique ground state $\ket{0}$ with $\hat{a}\ket{0}=0$ exists, and the
remaining eigenstates are
$\ket{n}=(\hat{a}^\dagger)^n/\sqrt{n!}\,\ket{0}$.

\begin{example}[Photons in a cavity]
A single mode of the electromagnetic field, treated in
\cref{ch:em-quantization}, has exactly this Hamiltonian with
$\op{N}$ counting photons. The same algebra describes phonons in a
crystal and, with the right choice of variables, the LC oscillator
of a superconducting circuit (\cref{ch:circuit-qed}).
\end{example}

\paragraph{Transition.}
The harmonic oscillator gives us an infinite ladder of evenly spaced
levels. The opposite extreme --- a system with only two levels --- is
the second cornerstone of quantum optics and forms the idealisation
of a qubit.

% --------------------------------------------------------------
\section{The two-level system and the Pauli algebra}
\label{sec:tls}

\begin{phenomenon}
Whenever a small subset of energy levels is well separated from
everything else, dynamics restricted to that subset can be modelled
by a Hamiltonian acting on a low-dimensional Hilbert space. The
simplest non-trivial case --- a single isolated transition --- is
the two-level system.
\end{phenomenon}

\begin{definition}[Pauli operators]
\label{def:pauli}
On $\hilbert=\C^2$ with basis $\{\ket{0},\ket{1}\}$, define
\begin{equation}
\hat\sigma_x = \begin{pmatrix}0&1\\1&0\end{pmatrix},\quad
\hat\sigma_y = \begin{pmatrix}0&-i\\i&0\end{pmatrix},\quad
\hat\sigma_z = \begin{pmatrix}1&0\\0&-1\end{pmatrix}.
\end{equation}
Together with $\id$ they form a basis of the Hermitian operators on
$\C^2$.
\end{definition}

The Pauli operators satisfy
\begin{equation}
\label{eq:pauli-algebra}
\hat\sigma_i\hat\sigma_j = \delta_{ij}\,\id + i\epsilon_{ijk}\hat\sigma_k,
\qquad
[\hat\sigma_i,\hat\sigma_j] = 2i\epsilon_{ijk}\hat\sigma_k,
\qquad
\{\hat\sigma_i,\hat\sigma_j\} = 2\delta_{ij}\,\id.
\end{equation}
We also use the raising and lowering operators
$\hat\sigma_\pm = (\hat\sigma_x\pm i\hat\sigma_y)/2$, with
$\hat\sigma_+\ket{1}=\ket{0}$, $\hat\sigma_-\ket{0}=\ket{1}$ in the
convention $\hat\sigma_z\ket{0}=+\ket{0}$ used throughout these notes.

\begin{statement}[Most general TLS Hamiltonian]
\label{stmt:tls-hamiltonian}
Up to an irrelevant additive constant, every Hermitian operator on
$\C^2$ has the form
\begin{equation}
\label{eq:tls-h}
\op{H} \;=\; \tfrac{\hbar}{2}\,\bm{\Omega}\cdot\hat{\bm{\sigma}},
\end{equation}
for some real vector $\bm{\Omega}=(\Omega_x,\Omega_y,\Omega_z)\in\R^3$
of angular frequencies. (We pull $\hbar$ out explicitly so that the
components of $\bm{\Omega}$ have units of $1/\text{time}$.)
\end{statement}

The eigenvalues of \cref{eq:tls-h} are $\pm\hbar|\bm{\Omega}|/2$, and the
eigenvectors are aligned with $\pm\bm{\Omega}/|\bm{\Omega}|$ in the geometric
sense made precise next.

\paragraph{Transition.}
The vector $\bm{\Omega}$ in \cref{eq:tls-h} suggests a geometric
picture. That picture is the Bloch sphere, and it turns the algebra
of a TLS into elementary geometry on $S^2$.

% --------------------------------------------------------------
\section{The Bloch sphere}
\label{sec:bloch}

\begin{statement}[Bloch representation]
\label{stmt:bloch}
Every state of a single qubit, pure or mixed, can be written as
\begin{equation}
\label{eq:bloch-rho}
\op{\rho} \;=\; \tfrac12\bigl(\id + \bm{r}\cdot\hat{\bm{\sigma}}\bigr),
\qquad
\bm{r}\in\R^3,\;|\bm{r}|\leq1,
\end{equation}
with $|\bm{r}|=1$ for pure states and $|\bm{r}|<1$ for mixed states.
\end{statement}

\begin{proof}
Any Hermitian operator on $\C^2$ expands in the basis
$\{\id,\hat\sigma_x,\hat\sigma_y,\hat\sigma_z\}$. Trace-one and
Hermiticity fix the coefficient of $\id$ and the reality of the
others; positivity $\op{\rho}\geq0$ requires both eigenvalues
$\tfrac12(1\pm|\bm{r}|)\geq0$, i.e.\ $|\bm{r}|\leq1$. Purity,
$\Tr\op{\rho}^2=\tfrac12(1+|\bm{r}|^2)=1$, then forces $|\bm{r}|=1$.
\end{proof}

In coordinates a pure state $\ket{\psi}=\cos(\theta/2)\ket{0}+
e^{i\varphi}\sin(\theta/2)\ket{1}$ has Bloch vector
$\bm{r}=(\sin\theta\cos\varphi,\sin\theta\sin\varphi,\cos\theta)$,
and the action of \cref{eq:tls-h} on $\bm{r}$ is precession around
$\bm{\Omega}$ at angular speed $|\bm{\Omega}|$. (The factor $\hbar$
in \cref{eq:tls-h} cancels with the $\hbar$ in
$e^{-i\op{H}t/\hbar}$.) Density matrices and the appearance of
$|\bm{r}|<1$ for mixed states are revisited in \cref{ch:composite}.

\begin{remark}
Single-qubit gates correspond to rotations of $\bm{r}$ about an axis
$\hat{\bm{n}}$ by an angle $\phi$, implemented by the unitary
\begin{equation}
\hat{R}_{\hat{\bm{n}}}(\phi) = \exp\!\bigl(-i\tfrac{\phi}{2}\,\hat{\bm{n}}\cdot\hat{\bm{\sigma}}\bigr)
= \cos(\phi/2)\,\id - i\sin(\phi/2)\,\hat{\bm{n}}\cdot\hat{\bm{\sigma}}.
\end{equation}
This is the Hamiltonian-engineering identity we use to think about
gate calibration in \cref{sec:rabi}.
\end{remark}

\paragraph{Transition.}
The static spectrum of a TLS is now fully characterised. The
question of how a TLS \emph{responds to a coherent drive} is what
turns this static picture into a tool for controlling qubits.

% --------------------------------------------------------------
\section{Rabi oscillations}
\label{sec:rabi}

\begin{phenomenon}
A two-level system irradiated near its transition frequency does
not absorb a single photon and stay in the excited state forever:
it oscillates coherently between $\ket{0}$ and $\ket{1}$ at a rate
controlled by the drive amplitude. The frequency of this oscillation
--- the Rabi frequency --- is the basic quantity any qubit
calibration aims to measure.
\end{phenomenon}

Consider the driven TLS Hamiltonian
\begin{equation}
\label{eq:driven-tls}
\op{H}(t) = \frac{\hbar\omega_q}{2}\hat\sigma_z
          + \hbar\Omega\cos(\omega_d t)\,\hat\sigma_x,
\end{equation}
with qubit (angular) frequency $\omega_q$, drive frequency $\omega_d$
and drive amplitude $\Omega$ (also an angular frequency).

\paragraph{Move to the rotating frame.}
Apply the unitary $\op{U}_0(t)=e^{-i\omega_d t\hat\sigma_z/2}$. The
transformed Hamiltonian, including the inertial term
$-i\hbar\,\op{U}_0^\dagger\dot{\op{U}}_0=-\hbar\omega_d\hat\sigma_z/2$, is
\begin{equation}
\op{H}_R(t) = \frac{\hbar\Delta}{2}\hat\sigma_z
            + \hbar\Omega\cos(\omega_d t)\,\op{U}_0^\dagger\hat\sigma_x\op{U}_0,
\qquad \Delta=\omega_q-\omega_d.
\end{equation}
Using $\op{U}_0^\dagger\hat\sigma_\pm\op{U}_0=e^{\pm i\omega_d t}\hat\sigma_\pm$
gives
$\op{U}_0^\dagger\hat\sigma_x\op{U}_0=
e^{i\omega_d t}\hat\sigma_+ + e^{-i\omega_d t}\hat\sigma_-$. Multiplying
by $\cos(\omega_d t)=\tfrac12(e^{i\omega_d t}+e^{-i\omega_d t})$ produces
four terms; two are time-independent and two oscillate at $\pm 2\omega_d$.
Dropping the latter (the rotating-wave approximation of \cref{sec:rwa}),
\begin{equation}
\label{eq:rabi-h-rwa}
\op{H}_R = \frac{\hbar\Delta}{2}\hat\sigma_z + \frac{\hbar\Omega}{2}\hat\sigma_x.
\end{equation}
This is exactly the form \cref{eq:tls-h} with
$\bm{\Omega}=(\Omega,0,\Delta)$, so the Bloch vector precesses about
$\bm{\Omega}$ at the \emph{generalised Rabi frequency}
\begin{equation}
\Omega_R = \sqrt{\Omega^2+\Delta^2}.
\end{equation}

\begin{theorem}[Rabi formula]
\label{thm:rabi}
For the driven TLS \cref{eq:driven-tls}, in the rotating-wave
approximation with the qubit initially in $\ket{0}$, the probability
of finding it in $\ket{1}$ at time $t$ is
\begin{equation}
\label{eq:rabi-formula}
P_{0\to 1}(t) \;=\; \frac{\Omega^2}{\Omega^2+\Delta^2}\,
\sin^2\!\Bigl(\frac{\Omega_R\,t}{2}\Bigr).
\end{equation}
\end{theorem}

\begin{proof}
The unitary generated by $\op{H}_R$ over a time $t$ is
\begin{equation}
e^{-i\op{H}_R t/\hbar}
= \cos(\Omega_R t/2)\,\id - i\sin(\Omega_R t/2)\,\hat{\bm{n}}\cdot\hat{\bm{\sigma}},
\end{equation}
with $\hat{\bm{n}}=(\Omega,0,\Delta)/\Omega_R$ (the explicit factors of
$\hbar$ in $\op{H}_R$ and in the Schr\"odinger evolution cancel). Acting
on $\ket{0}$ and projecting onto $\ket{1}$ gives \cref{eq:rabi-formula}.
\end{proof}

\begin{example}[Rabi calibration of a qubit]
On resonance $\Delta=0$ and $P_{0\to1}(t)=\sin^2(\Omega t/2)$. A
$\pi$-pulse, which inverts $\ket{0}\leftrightarrow\ket{1}$, requires
$\Omega t=\pi$. Sweeping the drive amplitude (or duration) and
fitting the resulting oscillation in $P_{0\to1}$ is the standard
``Rabi experiment'' used to calibrate the amplitude of single-qubit
gates in superconducting hardware (\cref{sec:experimental-basics}).
\end{example}

% --------------------------------------------------------------
\section*{Chapter summary and outlook}

We have solved the two systems that recur throughout the rest of the
text. The harmonic oscillator gave us the ladder operators, the Fock
basis, and the equally spaced spectrum that defines a linear
resonator. The two-level system gave us the Pauli algebra, the
geometric picture of the Bloch sphere, and Rabi oscillations under
a coherent drive.

So far our quantum systems have been completely isolated and
described by a single state vector. Real experiments are not isolated:
the qubit and resonator are coupled, the system is entangled with
its environment, and the experimenter has only partial information
about its state. The next chapter introduces the language --- tensor
products, density matrices, and the partial trace --- that is needed
to handle composite and mixed states without giving up any of the
machinery developed so far.
